\documentclass[11pt]{amsart}
\usepackage[margin=28mm]{geometry}
\usepackage{amssymb,tikz,url}
\usepackage[hidelinks]{hyperref}
\newtheorem{theorem}{Theorem}
\newtheorem{lemma}[theorem]{Lemma}
\theoremstyle{remark}
\newtheorem*{remark}{Remark}

\title{A Hamiltonian cubic graph with no cycle of length $n-1$}
\author{Vamshi Jandhyala}
\date{Preprint, 9 October 2026. Not peer reviewed.}

\begin{document}
\begin{abstract}
Gordon Royle asked whether some non-bipartite Hamiltonian cubic graph on $n$ vertices has no cycle of
length $n-1$, and reported that a computer found none on up to $24$ vertices. Replacing each vertex of
$K_4$ by a copy of $K_{2,3}$ gives one on $20$ vertices. The proof is a two-line count of cycle edges
in one block, and it applies to any Hamiltonian non-bipartite cubic graph in place of $K_4$. For the
$20$-vertex graph every statement is formally verified in Lean.
\end{abstract}
\maketitle

\section{The question}

A cycle of length $n$ in a graph on $n$ vertices is Hamiltonian. In a bipartite graph every cycle has
even length, so a Hamiltonian bipartite graph on $n$ vertices has no cycle of length $n-1$; the cube is
an example. In 2017 Royle asked on MathOverflow~\cite{mo} whether a cubic graph can be Hamiltonian,
not bipartite, and still have no cycle of length $n-1$, and reported that his computer found none on up
to $24$ vertices. The question has had no answer there.

\begin{theorem}\label{thm:main}
Let $G$ be the graph obtained from $K_4$ by replacing each vertex by a copy of $K_{2,3}$, the three
edges at that vertex being attached to the three vertices of the larger class. Then $G$ is a cubic
graph on $20$ vertices that is Hamiltonian, is not bipartite, and has no cycle of length $19$.
\end{theorem}

\section{The construction}

Start from any cubic graph $K$. Replace each vertex $v$ by a block $X_v$, a copy of $K_{2,3}$. Call
its two vertices of degree three the \emph{inner} vertices and its three vertices of degree two the
\emph{attachments}, and join the three attachments of $X_v$ to the three edges of $K$ at $v$, one each.
Every vertex now has degree three, and the graph has $5|V(K)|$ vertices. With $K=K_4$ this is the
graph of Theorem~\ref{thm:main}, drawn in Figure~\ref{fig:graph}.

\begin{figure}[ht]
\centering
\begin{tikzpicture}[scale=0.62]
\draw[black!35,line width=0.6pt] (-2.717,3.919) -- (-1.500,3.000);
\draw[orange!80!black,line width=2.2pt] (-2.717,3.919) -- (-1.939,1.939);
\draw[orange!80!black,line width=2.2pt] (-2.717,3.919) -- (-3.000,1.500);
\draw[orange!80!black,line width=2.2pt] (-3.919,2.717) -- (-1.500,3.000);
\draw[orange!80!black,line width=2.2pt] (-3.919,2.717) -- (-1.939,1.939);
\draw[black!35,line width=0.6pt] (-3.919,2.717) -- (-3.000,1.500);
\draw[orange!80!black,line width=2.2pt] (-1.500,3.000) -- (1.500,3.000);
\draw[black!35,line width=0.6pt] (-1.939,1.939) -- (1.939,-1.939);
\draw[orange!80!black,line width=2.2pt] (-3.000,1.500) -- (-3.000,-1.500);
\draw[black!35,line width=0.6pt] (3.919,2.717) -- (3.000,1.500);
\draw[orange!80!black,line width=2.2pt] (3.919,2.717) -- (1.939,1.939);
\draw[orange!80!black,line width=2.2pt] (3.919,2.717) -- (1.500,3.000);
\draw[orange!80!black,line width=2.2pt] (2.717,3.919) -- (3.000,1.500);
\draw[orange!80!black,line width=2.2pt] (2.717,3.919) -- (1.939,1.939);
\draw[black!35,line width=0.6pt] (2.717,3.919) -- (1.500,3.000);
\draw[orange!80!black,line width=2.2pt] (3.000,1.500) -- (3.000,-1.500);
\draw[black!35,line width=0.6pt] (1.939,1.939) -- (-1.939,-1.939);
\draw[black!35,line width=0.6pt] (2.717,-3.919) -- (1.500,-3.000);
\draw[orange!80!black,line width=2.2pt] (2.717,-3.919) -- (1.939,-1.939);
\draw[orange!80!black,line width=2.2pt] (2.717,-3.919) -- (3.000,-1.500);
\draw[orange!80!black,line width=2.2pt] (3.919,-2.717) -- (1.500,-3.000);
\draw[orange!80!black,line width=2.2pt] (3.919,-2.717) -- (1.939,-1.939);
\draw[black!35,line width=0.6pt] (3.919,-2.717) -- (3.000,-1.500);
\draw[orange!80!black,line width=2.2pt] (1.500,-3.000) -- (-1.500,-3.000);
\draw[black!35,line width=0.6pt] (-3.919,-2.717) -- (-3.000,-1.500);
\draw[orange!80!black,line width=2.2pt] (-3.919,-2.717) -- (-1.939,-1.939);
\draw[orange!80!black,line width=2.2pt] (-3.919,-2.717) -- (-1.500,-3.000);
\draw[orange!80!black,line width=2.2pt] (-2.717,-3.919) -- (-3.000,-1.500);
\draw[orange!80!black,line width=2.2pt] (-2.717,-3.919) -- (-1.939,-1.939);
\draw[black!35,line width=0.6pt] (-2.717,-3.919) -- (-1.500,-3.000);
\filldraw[fill=white,draw=black] (-2.717,3.919) circle (0.13);
\filldraw[fill=white,draw=black] (-3.919,2.717) circle (0.13);
\filldraw[fill=black,draw=black] (-1.500,3.000) circle (0.13);
\filldraw[fill=black,draw=black] (-1.939,1.939) circle (0.13);
\filldraw[fill=black,draw=black] (-3.000,1.500) circle (0.13);
\filldraw[fill=white,draw=black] (3.919,2.717) circle (0.13);
\filldraw[fill=white,draw=black] (2.717,3.919) circle (0.13);
\filldraw[fill=black,draw=black] (3.000,1.500) circle (0.13);
\filldraw[fill=black,draw=black] (1.939,1.939) circle (0.13);
\filldraw[fill=black,draw=black] (1.500,3.000) circle (0.13);
\filldraw[fill=white,draw=black] (2.717,-3.919) circle (0.13);
\filldraw[fill=white,draw=black] (3.919,-2.717) circle (0.13);
\filldraw[fill=black,draw=black] (1.500,-3.000) circle (0.13);
\filldraw[fill=black,draw=black] (1.939,-1.939) circle (0.13);
\filldraw[fill=black,draw=black] (3.000,-1.500) circle (0.13);
\filldraw[fill=white,draw=black] (-3.919,-2.717) circle (0.13);
\filldraw[fill=white,draw=black] (-2.717,-3.919) circle (0.13);
\filldraw[fill=black,draw=black] (-3.000,-1.500) circle (0.13);
\filldraw[fill=black,draw=black] (-1.939,-1.939) circle (0.13);
\filldraw[fill=black,draw=black] (-1.500,-3.000) circle (0.13);
\end{tikzpicture}
\qquad\qquad
\begin{tikzpicture}[scale=0.62]
\draw[blue!70!black,line width=2.2pt] (-2.717,3.919) -- (-1.500,3.000);
\draw[blue!70!black,line width=2.2pt] (-2.717,3.919) -- (-1.939,1.939);
\draw[black!35,line width=0.6pt] (-2.717,3.919) -- (-3.000,1.500);
\draw[black!35,line width=0.6pt] (-3.919,2.717) -- (-1.500,3.000);
\draw[black!35,line width=0.6pt] (-3.919,2.717) -- (-1.939,1.939);
\draw[black!35,line width=0.6pt] (-3.919,2.717) -- (-3.000,1.500);
\draw[blue!70!black,line width=2.2pt] (-1.500,3.000) -- (1.500,3.000);
\draw[blue!70!black,line width=2.2pt] (-1.939,1.939) -- (1.939,-1.939);
\draw[black!35,line width=0.6pt] (-3.000,1.500) -- (-3.000,-1.500);
\draw[blue!70!black,line width=2.2pt] (3.919,2.717) -- (3.000,1.500);
\draw[black!35,line width=0.6pt] (3.919,2.717) -- (1.939,1.939);
\draw[blue!70!black,line width=2.2pt] (3.919,2.717) -- (1.500,3.000);
\draw[black!35,line width=0.6pt] (2.717,3.919) -- (3.000,1.500);
\draw[black!35,line width=0.6pt] (2.717,3.919) -- (1.939,1.939);
\draw[black!35,line width=0.6pt] (2.717,3.919) -- (1.500,3.000);
\draw[blue!70!black,line width=2.2pt] (3.000,1.500) -- (3.000,-1.500);
\draw[black!35,line width=0.6pt] (1.939,1.939) -- (-1.939,-1.939);
\draw[black!35,line width=0.6pt] (2.717,-3.919) -- (1.500,-3.000);
\draw[blue!70!black,line width=2.2pt] (2.717,-3.919) -- (1.939,-1.939);
\draw[blue!70!black,line width=2.2pt] (2.717,-3.919) -- (3.000,-1.500);
\draw[black!35,line width=0.6pt] (3.919,-2.717) -- (1.500,-3.000);
\draw[black!35,line width=0.6pt] (3.919,-2.717) -- (1.939,-1.939);
\draw[black!35,line width=0.6pt] (3.919,-2.717) -- (3.000,-1.500);
\draw[black!35,line width=0.6pt] (1.500,-3.000) -- (-1.500,-3.000);
\draw[black!35,line width=0.6pt] (-3.919,-2.717) -- (-3.000,-1.500);
\draw[black!35,line width=0.6pt] (-3.919,-2.717) -- (-1.939,-1.939);
\draw[black!35,line width=0.6pt] (-3.919,-2.717) -- (-1.500,-3.000);
\draw[black!35,line width=0.6pt] (-2.717,-3.919) -- (-3.000,-1.500);
\draw[black!35,line width=0.6pt] (-2.717,-3.919) -- (-1.939,-1.939);
\draw[black!35,line width=0.6pt] (-2.717,-3.919) -- (-1.500,-3.000);
\filldraw[fill=white,draw=black] (-2.717,3.919) circle (0.13);
\filldraw[fill=white,draw=black] (-3.919,2.717) circle (0.13);
\filldraw[fill=black,draw=black] (-1.500,3.000) circle (0.13);
\filldraw[fill=black,draw=black] (-1.939,1.939) circle (0.13);
\filldraw[fill=black,draw=black] (-3.000,1.500) circle (0.13);
\filldraw[fill=white,draw=black] (3.919,2.717) circle (0.13);
\filldraw[fill=white,draw=black] (2.717,3.919) circle (0.13);
\filldraw[fill=black,draw=black] (3.000,1.500) circle (0.13);
\filldraw[fill=black,draw=black] (1.939,1.939) circle (0.13);
\filldraw[fill=black,draw=black] (1.500,3.000) circle (0.13);
\filldraw[fill=white,draw=black] (2.717,-3.919) circle (0.13);
\filldraw[fill=white,draw=black] (3.919,-2.717) circle (0.13);
\filldraw[fill=black,draw=black] (1.500,-3.000) circle (0.13);
\filldraw[fill=black,draw=black] (1.939,-1.939) circle (0.13);
\filldraw[fill=black,draw=black] (3.000,-1.500) circle (0.13);
\filldraw[fill=white,draw=black] (-3.919,-2.717) circle (0.13);
\filldraw[fill=white,draw=black] (-2.717,-3.919) circle (0.13);
\filldraw[fill=black,draw=black] (-3.000,-1.500) circle (0.13);
\filldraw[fill=black,draw=black] (-1.939,-1.939) circle (0.13);
\filldraw[fill=black,draw=black] (-1.500,-3.000) circle (0.13);
\end{tikzpicture}

\caption{$K_4$ with each vertex replaced by $K_{2,3}$: inner vertices white, attachments black.
Left, a Hamiltonian cycle. Right, a cycle of length nine.}\label{fig:graph}
\end{figure}

\emph{Hamiltonian.} Any two attachments $p$, $p'$ of a block are the ends of a path through all five
of its vertices: $p$, an inner vertex, the third attachment, the other inner vertex, $p'$. Following a
Hamiltonian cycle of $K$ and crossing each block by such a path gives a Hamiltonian cycle of the new
graph (Figure~\ref{fig:graph}, left).

\emph{Not bipartite.} Two attachments of a block are also joined by a path of length two through an
inner vertex. Following a triangle of $K$ this way gives a cycle of length $3\cdot2+3=9$
(Figure~\ref{fig:graph}, right).

\section{No cycle of length \texorpdfstring{$n-1$}{n-1}}

\begin{lemma}\label{lem:count}
If $K$ is any cubic graph, the graph built from it has no cycle of length $n-1$.
\end{lemma}

\begin{proof}
Suppose a cycle $C$ passes through every vertex except one, say $m$, which lies in the block $X$.
The other four vertices of $X$ lie on $C$ and each has two edges of $C$. Count these eight edge-ends.
An edge of $C$ with both ends in $X$ is counted twice and an edge of $C$ leaving $X$ once, so if $C$
has $i$ edges inside $X$ and $e$ edges leaving it,
\[
2i+e=8.
\]
Every edge inside $X$ joins an inner vertex to an attachment, and an inner vertex has no neighbours
outside $X$. So each inner vertex on $C$ contributes both of its $C$-edges to $i$, and no edge inside
$X$ is counted from two inner vertices: $i$ is twice the number of inner vertices on $C$.

If $m$ is an inner vertex, one inner vertex is on $C$, so $i=2$ and $e=4$. But only three edges leave
$X$, one at each attachment. If $m$ is an attachment, both inner vertices are on $C$, so $i=4$ and
$e=0$. Then $C$ never leaves $X$ and has at most five vertices, while $n-1\ge 19$. Either way we reach
a contradiction.
\end{proof}

Theorem~\ref{thm:main} follows. The same argument applies to every bipartite block attached by three
edges whose larger class carries the attachments; and replacing $K_4$ by any Hamiltonian cubic graph
that is not bipartite, such as the triangular prism, gives further examples, on $30$ vertices for the
prism.

\begin{remark}
The theorem disagrees with the search reported in the question, and we do not know how that search
was restricted. Our graph has girth four and three-edge cuts that separate a block from the rest, so
it lies outside the classes of cubic graphs most often searched, such as those of girth five or those
that are cyclically $4$-edge-connected. Whether an example exists in those classes, and the least
number of vertices of any example, we leave open.
\end{remark}

\section*{Verification and use of AI tools}

For the $20$-vertex graph every statement of Theorem~\ref{thm:main} is formally verified in Lean~4
with Mathlib (folder \path{lean}, standard axioms only). The theorem about cycles is proved for every
cycle, in Mathlib's sense, by the count of Lemma~\ref{lem:count}; the finite facts about the fixed graph
(its degrees, its neighbourhoods and the two cycles of Figure~\ref{fig:graph}) are checked by
evaluation. Six deliberately wrong variants of the formal statements fail to compile. Independently,
the program \path{verify/check_cycles.py} builds the graph separately and confirms by exhaustive search
that deleting any one vertex leaves a non-Hamiltonian graph, with the cube and the prism as controls; it
finds the same for the $30$-vertex example. The author used an AI tool (Claude, Anthropic) to find the
example and write the formal proof, and is responsible for the content.

\begin{thebibliography}{9}
\bibitem{mo} G.~Royle, Is there a non-bipartite hamiltonian cubic graph on $n$ vertices with no
$(n-1)$-cycle?, MathOverflow question 263706 (2017), \url{https://mathoverflow.net/q/263706},
accessed 9 October 2026.
\end{thebibliography}
\end{document}
